\Askhsh{37104}{4}
\begin{ylh}{1}
    \item 4.6. Άθροισμα γωνιών τριγώνου
    \item 5.8. Το ορθόκεντρο τριγώνου
    \item 5.9. Μια ιδιότητα του ορθογώνιου τριγώνου.
\end{ylh}
% ===================================================================
%  Αρχείο: 37104.tex (Fragment)
%  Αυτόματη μετατροπή μέσω doc2latex.py
% ===================================================================

ΘΕΜΑ 4

Δίνεται ορθογώνιο τρίγωνο ΑΒΓ ($\widehat{Α} = 90^{ο}$) και $\widehat{\Gamma} = 30^{ο}$ με Μ και Ν τα μέσα των πλευρών ΒΓ και ΑΒ αντίστοιχα. Έστω ότι η μεσοκάθετος της πλευράς ΒΓ τέμνει την ΑΓ στο σημείο Ε.

\begin{alist}
\item Να αποδείξετε ότι:

i) η ΒΕ είναι διχοτόμος της γωνίας $\widehat{Β}$. \monades{6}

ii) $ΑΕ = \ \frac{\Gamma Ε}{2}$. \monades{6}

iii) η ΒΕ είναι μεσοκάθετος της διαμέσου ΑΜ. \monades{7}

\item Αν ΑΔ είναι το ύψος του τριγώνου ΑΒΓ που τέμνει την ΒΕ στο Η, να αποδείξετε ότι τα σημεία Μ, Η και Ν είναι συνευθειακά. \monades{6}

\includesvg[width=5.50889in,height=2.97917in]{/home/spyros/Μαθηματικά/Trapeza_Thematwn/A_Lykeiou/Geometry/extracted/37104/media/image2.svg}
\end{alist}
